\section{Analysis}\label{sec:analysis}

\begin{figure*}[!t]
  \centering
  \includegraphics[width=\textwidth]{fig3_gate_shift.pdf}
  \caption{(a)~Gate turn conditioning: per-layer gate values at turn~1 vs turn~21. In 27/28 layers, gates approach 1.0 at late turns, inverting the design intent ($p = 3.73 \times 10^{-8}$). (b)~Distribution shift ($N$\,=\,5 pilot): the TF-AR GAR gap persists across all five variants. Pilot gaps (20--34$\times$) exceed final $N$\,=\,50 estimates (8--17$\times$; Table~\ref{tab:five-way}) due to smaller sample size.}
  \label{fig:gate-shift}
\end{figure*}

Each subsection kills one alternative hypothesis via diagnostic exclusion: \S\ref{subsec:gate-conditioning}~turn conditioning, \S\ref{subsec:pcgrad}~gradient direction conflict, leaving exposure bias as the primary remaining explanation consistent with current evidence (\S\ref{subsec:dist-shift}). \S\ref{subsec:mechanism-specificity} confirms selectivity; \S\ref{subsec:cross-model} confirms cross-architecture generalization.

\subsection{Gate Turn Conditioning}\label{subsec:gate-conditioning}

We compare gate values at turn~1 (early) and turn~21 (late) across all 28 layers. The gates invert the design intent: in 27/28 layers, gate values are \emph{lower} at turn~1 than turn~21 (Wilcoxon $p = 3.73 \times 10^{-8}$, $r = 0.975$); e.g., layer~23 shifts from 0.769 to 0.979. NTP pressure pushes gates toward 1.0, a convergence pattern reminiscent of attention collapse in deep transformers \citep{dong2021attention}, yielding inference-time values of 0.93--0.98 (2--7\% attenuation; Table~\ref{tab:five-way}; full gate profiles in Appendix~\ref{sec:appendix-gate-profile}; training dynamics in Appendix~\ref{sec:appendix-training-ext}). The Uniform Gate learns comparable per-layer variation (L3\,=\,0.777, L27\,=\,1.000) despite lacking turn conditioning; the structure thus arises from LoRA$+$GAR pressure (Figure~\ref{fig:gate-shift}a). Does the NTP-GAR gradient conflict prevent the intended behavior?

\subsection{PCGrad Gradient Direction Conflict Exclusion}\label{subsec:pcgrad}

NTP and GAR gradients oppose in 68--79\% of layers (cosine similarity $-$0.089$\to$$-$0.042 across epochs), with NTP magnitude dominance growing from 7.5$\times$ to 14.7$\times$ (full gradient statistics in Appendix~\ref{sec:appendix-gradient}). If this direction conflict prevents GAR convergence, eliminating it via PCGrad \citep{yu2020pcgrad}, which projects conflicting gradients onto the other's normal plane, should prevent the collapse. To our knowledge, this is the first application of PCGrad as a \emph{diagnostic} tool rather than an optimization method.

PCGrad shifts cosine similarity to oscillate around zero, yet TF metrics are unchanged. AR GAR \emph{worsens}: 0.064$\to$0.030 ($p$\,=\,0.002; Figure~\ref{fig:training-pcgrad}b in Appendix~\ref{sec:appendix-gradient}). By \emph{modus tollens}: removing gradient direction conflict does not prevent collapse; therefore it is not the primary cause. PCGrad eliminates direction conflict but does not address magnitude imbalance (NTP dominates GAR by 7.5--14.7$\times$); the worsening under PCGrad may partly reflect projection side-effects rather than confirming that direction conflict is entirely absent as a factor. Cross-model analysis confirms this pattern: Llama-3.1-8B gradients are near-orthogonal (cos\,=\,+0.005) with 25.7$\times$ NTP magnitude dominance (Table~\ref{tab:gradient-crossmodel}).

\subsection{Distribution Shift Analysis}\label{subsec:dist-shift}

The TF-AR GAR gap persists across all variants (Table~\ref{tab:five-way}). GAR-Dongre provides a control: it remains uniformly low (0.01--0.02) with no distribution gap, evidence that the shift selectively destroys concentrated last-token attention. The mechanism is exposure bias \citep{bengio2015scheduled,zhang2019bridging,arora2022exposure}: during training the model conditions on ground-truth prefixes; at inference, on its own outputs. GAR optimizes a distributional property with no autoregressive counterpart, explaining why TF GAR-ours remains near $\tau$\,=\,0.5 while AR GAR collapses (77--87\% per-turn decay; Figure~\ref{fig:attention-heatmap}; per-configuration AR trajectories in Appendix~\ref{sec:appendix-ar-gar-dynamics}; scaling analysis in Appendix~\ref{sec:appendix-eb-scaling}).

\subsection{Mechanism Specificity}\label{subsec:mechanism-specificity}

We test specificity with Turn-level Residual Coherence (TRC), the cosine similarity between adjacent-turn hidden states. TF/AR TRC is 0.334/0.332 for ThreadGate (gap $+$0.002) and 0.349/0.411 for LoRA-only (gap $-$0.062, reversed); exposure bias thus disrupts a specific routing behavior while leaving representations intact. Cross-model evaluation confirms: GAR-Dongre shows negligible shift on both Llama (1.002$\times$) and Qwen (1.28$\times$) vs.\ the order-of-magnitude GAR-ours collapse.

\begin{figure}[t]
  \centering
  \includegraphics[width=\columnwidth]{fig_ss_multiseed.pdf}
  \caption{Multi-seed AR GAR under scheduled sampling ($\beta$\,=\,0 and $\beta$\,=\,5, three seeds). Steering nearly eliminates seed variance (std: 0.047$\to$0.001).}
  \label{fig:ss-multiseed}
\end{figure}

\subsection{Cross-Model Replication}\label{subsec:cross-model}

We replicate on Llama-3.1-8B-Instruct \citep{grattafiori2024llama3} (Base, ThreadGate, LoRA-only; $N$\,=\,20, 4 turns). TF GAR-ours is 0.484 (TG) / 0.481 (LoRA-only); gating contributes zero benefit. The AR distribution is heavy-tailed (coefficient of variation, CoV\,$>$\,100\%); we report medians: 0.018 (TG and LoRA-only), 0.009 (Base), yielding gaps of 27$\times$ and 26$\times$ (Figure~\ref{fig:crossmodel}a in Appendix~\ref{sec:appendix-llama}; per-turn values in Appendix~\ref{sec:appendix-perturn-detail}). Gate profiles differ qualitatively (Figure~\ref{fig:crossmodel}b): Qwen attenuates early layers (L1\,=\,0.793) while Llama dips at late layers (L25\,=\,0.719), an architecture-dependent pattern, but convergence toward identity is universal (Appendix~\ref{sec:appendix-viz}). The gap is substantial on both architectures, consistent with exposure bias as architecture-independent.

Scheduled sampling confirms cross-model generalization: Llama SS compresses the gap from ${\sim}$27$\times$ to 2.9$\times$ (TF GAR 0.477, AR GAR 0.166), weaker than Qwen's ${\sim}$1.3$\times$ likely due to fewer training epochs (2 vs.\ 3). The phase transition is architecture-dependent: $\beta$\,=\,3 triggers recovery on Qwen SS (0.593) but not Llama SS (0.173); Llama requires $\beta$\,=\,5 (GAR 0.791, exceeding Qwen SS 0.717; Appendix~\ref{sec:appendix-llama-ss}). Cross-scale validation on Qwen-2.5-1.5B confirms the same pattern (79.8\% GAR decay, $\beta$-steering recovery to 0.715; Appendix~\ref{sec:appendix-1.5b}).

